\documentclass{article}
\usepackage{amsmath,amssymb,amsthm}
\usepackage{enumitem}
\usepackage{hyperref}
\usepackage{geometry}
\geometry{margin=1in}
\newtheorem{theorem}{Theorem}
\newtheorem{lemma}{Lemma}
\newtheorem{assumption}{Assumption}
\newcommand{\E}{\mathbb{E}}
\newcommand{\norm}[1]{\left\|#1\right\|}
\newcommand{\inner}[2]{\langle #1,#2\rangle}
\begin{document}

\section*{Algorithm Name}
\textbf{Nesterov‑type Accelerated Gradient with Polynomial Momentum}\\
(denoted \texttt{AG\_Poly})

\section*{Mathematical Formulation}
Let $f:\mathbb{R}^{d}\rightarrow\mathbb{R}$ be differentiable.  
Initialize $x_{0}=x_{-1}\in\mathbb{R}^{d}$ and fix a stepsize $\alpha>0$.  
For $k=0,1,2,\dots$ perform
\[
\begin{aligned}
\beta_{k}&=\frac{k-1}{k+2},\qquad (\beta_{0}:=0)\\[4pt]
y_{k}&=x_{k}+\beta_{k}\bigl(x_{k}-x_{k-1}\bigr),\\[4pt]
x_{k+1}&=y_{k}-\alpha\,\nabla f\!\bigl(y_{k}\bigr).
\end{aligned}
\tag{1}
\]

\section*{Pseudocode}
\begin{verbatim}
Input: stepsize α>0, initial point x0, maxIter K
Set x_{-1} ← x0
for k = 0,…,K-1 do
    if k==0 then β_k ← 0 else β_k ← (k-1)/(k+2)
    y_k ← x_k + β_k * (x_k - x_{k-1})
    g_k ← ∇f(y_k)
    x_{k+1} ← y_k - α * g_k
end for
return {x_k}_{k=0}^K
\end{verbatim}

\section*{Dry Run Example}
Consider the scalar quadratic 
\[
f(w)=(w-3)^{2},\qquad \nabla f(w)=2(w-3).
\]
Choose $\alpha=0.1$, $x_{0}=0$, and set $x_{-1}=x_{0}=0$.

\begin{center}
\begin{tabular}{c|c|c|c|c}
$k$ & $\beta_{k}$ & $y_{k}$ & $x_{k+1}=y_{k}-\alpha\nabla f(y_{k})$ & $f(x_{k+1})$\\\hline
0 & $0$ & $0+0(0-0)=0$ & $0-0.1\cdot(-6)=0.6$ & $(0.6-3)^{2}=5.76$\\
1 & $\frac{0}{3}=0$ & $0.6+0(0.6-0)=0.6$ & $0.6-0.1\cdot2(0.6-3)=0.6-0.1\cdot(-4.8)=1.08$ & $(1.08-3)^{2}=3.6864$\\
2 & $\frac{1}{4}=0.25$ & $1.08+0.25(1.08-0.6)=1.08+0.25(0.48)=1.2$ &
$1.2-0.1\cdot2(1.2-3)=1.2-0.1\cdot(-3.6)=1.56$ & $(1.56-3)^{2}=2.0736$
\end{tabular}
\end{center}
The objective values decrease from $9$ (at $w=0$) to $2.07$ after three iterations.

\newpage
\section*{Assumptions}
\begin{assumption}[Smoothness]\label{as:smooth}
$f$ is $L$‑smooth, i.e.
\[
\forall u,v\in\mathbb{R}^{d}:\quad 
f(u)\le f(v)+\inner{\nabla f(v)}{u-v}+\frac{L}{2}\norm{u-v}^{2}.
\tag{2}
\]
\end{assumption}
\begin{assumption}[Lower Boundedness]\label{as:lb}
There exists $f^{\star}>-\infty$ such that $f(x)\ge f^{\star}$ for all $x$.
\end{assumption}
\begin{assumption}[Stepsize]\label{as:step}
The stepsize satisfies $0<\alpha\le \frac{1}{L}$.
\end{assumption}

\section*{Main Convergence Theorem}
\begin{theorem}[Non‑convex convergence of \texttt{AG\_Poly}]
\label{thm:main}
Under Assumptions \ref{as:smooth}–\ref{as:step}, the iterates generated by (1) satisfy for any $K\ge 0$
\[
\min_{0\le k\le K}\norm{\nabla f(y_{k})}^{2}
\;\le\;
\frac{2L\bigl(f(x_{0})-f^{\star}\bigr)}{\alpha(K+1)}.
\tag{3}
\]
Consequently,
\[
\frac{1}{K+1}\sum_{k=0}^{K}\norm{\nabla f(y_{k})}^{2}
\;\le\;
\frac{2L\bigl(f(x_{0})-f^{\star}\bigr)}{\alpha(K+1)}.
\tag{4}
\]
Thus the algorithm attains an $O(1/K)$ rate in the squared gradient norm.
\end{theorem}

\section*{Theoretical Properties}
\begin{itemize}[leftmargin=*]
\item \textbf{Stability.} The polynomial schedule $\beta_{k}=(k-1)/(k+2)$ satisfies $0\le\beta_{k}<1$ for all $k\ge0$, guaranteeing that the momentum term never overshoots the previous direction.
\item \textbf{Hyper‑parameter sensitivity.} The only free scalar is $\alpha$, which must respect $\alpha\le1/L$. Within this interval the bound (3) is monotone decreasing in $\alpha$, i.e. larger $\alpha$ (up to $1/L$) yields a tighter guarantee.
\item \textbf{Comparison to baselines.}
\begin{enumerate}
\item Plain Gradient Descent with stepsize $\alpha\le1/L$ enjoys the same $O(1/K)$ bound on $\norm{\nabla f(x_{k})}^{2}$ but without any acceleration on the function values.
\item Classical Nesterov acceleration with constant $\beta\in[0,1)$ does \emph{not} guarantee a $O(1/K)$ bound for general non‑convex $f$ unless additional curvature conditions hold. The polynomial schedule restores the $O(1/K)$ guarantee while still providing empirical acceleration.
\end{enumerate}
\end{itemize}

\section*{Preliminaries (Definitions and Lemmas)}
\begin{lemma}[Descent Lemma]\label{lem:descent}
If $f$ is $L$‑smooth and $\alpha\le 1/L$, then for any $z$,
\[
f\bigl(z-\alpha\nabla f(z)\bigr)\le f(z)-\frac{\alpha}{2}\norm{\nabla f(z)}^{2}.
\tag{5}
\]
\end{lemma}
\begin{proof}
From \eqref{as:smooth} with $u=z-\alpha\nabla f(z)$ and $v=z$,
\[
f(z-\alpha\nabla f(z))\le f(z)-\alpha\norm{\nabla f(z)}^{2}
+\frac{L\alpha^{2}}{2}\norm{\nabla f(z)}^{2}.
\]
Because $\alpha\le1/L$, $L\alpha^{2}\le\alpha$, yielding \eqref{eq:descent}.
\end{proof}

\begin{lemma}[Momentum relation]\label{lem:beta}
For the schedule $\beta_{k}=\frac{k-1}{k+2}$ one has
\[
\frac{1-\beta_{k}}{\alpha}
=
\frac{3}{\alpha(k+2)}.
\tag{6}
\]
\end{lemma}
\begin{proof}
Simple algebra:
\[
1-\beta_{k}=1-\frac{k-1}{k+2}=\frac{3}{k+2},
\]
hence \eqref{6}.
\end{proof}

\section*{Main Proof (Step‑by‑Step Derivation)}
We prove Theorem \ref{thm:main}. All derivations are numbered for easy reference.

\noindent\textbf{Step 1.}  Write the update in terms of $x_{k}$ and $x_{k-1}$:
\[
x_{k+1}=x_{k}+\beta_{k}(x_{k}-x_{k-1})-\alpha\nabla f\!\bigl(y_{k}\bigr),
\qquad y_{k}=x_{k}+\beta_{k}(x_{k}-x_{k-1}).
\tag{7}
\]

\noindent\textbf{Step 2.}  Define the auxiliary potential
\[
\Phi_{k}:=f(x_{k})+\tfrac{c_{k}}{2\alpha}\norm{x_{k}-x_{k-1}}^{2},
\tag{8}
\]
with coefficients $c_{k}:= \frac{3}{k+2}$ (motivated by Lemma \ref{lem:beta}).

\noindent\textbf{Step 3.}  Compute $\Phi_{k+1}-\Phi_{k}$.
Using \eqref{eq:descent} on $y_{k}$ we obtain
\[
f(x_{k+1})\le f(y_{k})-\frac{\alpha}{2}\norm{\nabla f(y_{k})}^{2}.
\tag{9}
\]
Next, by smoothness of $f$,
\[
f(y_{k})\le f(x_{k})+\inner{\nabla f(x_{k})}{y_{k}-x_{k}}
+\frac{L}{2}\norm{y_{k}-x_{k}}^{2}.
\tag{10}
\]

\noindent\textbf{Step 4.}  Express $y_{k}-x_{k}= \beta_{k}(x_{k}-x_{k-1})$. Plugging into \eqref{10} gives
\[
f(y_{k})\le f(x_{k})+\beta_{k}\inner{\nabla f(x_{k})}{x_{k}-x_{k-1}}
+\frac{L\beta_{k}^{2}}{2}\norm{x_{k}-x_{k-1}}^{2}.
\tag{11}
\]

\noindent\textbf{Step 5.}  Combine \eqref{9} and \eqref{11}:
\[
f(x_{k+1})\le f(x_{k})+\beta_{k}\inner{\nabla f(x_{k})}{x_{k}-x_{k-1}}
+\frac{L\beta_{k}^{2}}{2}\norm{x_{k}-x_{k-1}}^{2}
-\frac{\alpha}{2}\norm{\nabla f(y_{k})}^{2}.
\tag{12}
\]

\noindent\textbf{Step 6.}  Add the quadratic term from $\Phi_{k+1}$:
\[
\Phi_{k+1}=f(x_{k+1})+\frac{c_{k+1}}{2\alpha}\norm{x_{k+1}-x_{k}}^{2}.
\tag{13}
\]
From the update rule,
\[
x_{k+1}-x_{k}= \beta_{k}(x_{k}-x_{k-1})-\alpha\nabla f(y_{k}).
\tag{14}
\]
Hence
\[
\begin{aligned}
\norm{x_{k+1}-x_{k}}^{2}
&= \beta_{k}^{2}\norm{x_{k}-x_{k-1}}^{2}
+\alpha^{2}\norm{\nabla f(y_{k})}^{2}
-2\alpha\beta_{k}\inner{\nabla f(y_{k})}{x_{k}-x_{k-1}}.
\end{aligned}
\tag{15}
\]

\noindent\textbf{Step 7.}  Substitute \eqref{15} into \eqref{13} and combine with \eqref{12}:
\[
\begin{aligned}
\Phi_{k+1}-\Phi_{k}
&\le
\beta_{k}\inner{\nabla f(x_{k})}{x_{k}-x_{k-1}}
+\frac{L\beta_{k}^{2}}{2}\norm{x_{k}-x_{k-1}}^{2}
-\frac{\alpha}{2}\norm{\nabla f(y_{k})}^{2}\\
&\quad
+\frac{c_{k+1}}{2\alpha}\Bigl[
\beta_{k}^{2}\norm{x_{k}-x_{k-1}}^{2}
+\alpha^{2}\norm{\nabla f(y_{k})}^{2}
-2\alpha\beta_{k}\inner{\nabla f(y_{k})}{x_{k}-x_{k-1}}
\Bigr].
\end{aligned}
\tag{16}
\]

\noindent\textbf{Step 8.}  Rearrange terms grouping the inner products and the squared norms.
First, the inner‑product part:
\[
\beta_{k}\inner{\nabla f(x_{k})}{x_{k}-x_{k-1}}
-\frac{c_{k+1}\beta_{k}}{\alpha}\inner{\nabla f(y_{k})}{x_{k}-x_{k-1}}.
\tag{17}
\]
Second, the $\norm{x_{k}-x_{k-1}}^{2}$ part:
\[
\frac{L\beta_{k}^{2}}{2}\norm{x_{k}-x_{k-1}}^{2}
+\frac{c_{k+1}\beta_{k}^{2}}{2\alpha}\norm{x_{k}-x_{k-1}}^{2}.
\tag{18}
\]
Third, the $\norm{\nabla f(y_{k})}^{2}$ part:
\[
-\frac{\alpha}{2}\norm{\nabla f(y_{k})}^{2}
+\frac{c_{k+1}\alpha}{2}\norm{\nabla f(y_{k})}^{2}
=
\frac{\alpha}{2}\bigl(c_{k+1}-1\bigr)\norm{\nabla f(y_{k})}^{2}.
\tag{19}
\]

\noindent\textbf{Step 9.}  Choose $c_{k}$ exactly as in \eqref{6}, i.e. $c_{k}= \frac{3}{k+2}$.  
Then one can verify
\[
c_{k+1}= \frac{3}{k+3},\qquad
c_{k+1}\beta_{k}= \frac{3}{k+3}\cdot\frac{k-1}{k+2}
= \frac{3(k-1)}{(k+2)(k+3)}.
\tag{20}
\]
Moreover,
\[
\frac{c_{k+1}\beta_{k}}{\alpha}
=
\frac{1-\beta_{k}}{\alpha}
\quad\text{by Lemma \ref{lem:beta}}.
\tag{21}
\]

\noindent\textbf{Step 10.}  Replace the inner‑product term \eqref{17} using \eqref{21}:
\[
\beta_{k}\inner{\nabla f(x_{k})}{x_{k}-x_{k-1}}
-\frac{1-\beta_{k}}{\alpha}\inner{\nabla f(y_{k})}{x_{k}-x_{k-1}}.
\tag{22}
\]
Add and subtract $\beta_{k}\inner{\nabla f(y_{k})}{x_{k}-x_{k-1}}$ to obtain
\[
\beta_{k}\inner{\nabla f(x_{k})-\nabla f(y_{k})}{x_{k}-x_{k-1}}
+\bigl(\beta_{k}-(1-\beta_{k})\bigr)\frac{1}{\alpha}
\inner{\nabla f(y_{k})}{x_{k}-x_{k-1}}.
\tag{23}
\]

\noindent\textbf{Step 11.}  By Cauchy–Schwarz and $L$‑smoothness,
\[
\norm{\nabla f(x_{k})-\nabla f(y_{k})}\le L\norm{x_{k}-y_{k}}
= L\beta_{k}\norm{x_{k}-x_{k-1}}.
\tag{24}
\]
Thus the first term of \eqref{23} is bounded by
\[
\beta_{k}L\beta_{k}\norm{x_{k}-x_{k-1}}^{2}
= L\beta_{k}^{2}\norm{x_{k}-x_{k-1}}^{2}.
\tag{25}
\]

\noindent\textbf{Step 12.}  Observe that $\beta_{k}-(1-\beta_{k})=2\beta_{k}-1 =\frac{2(k-1)}{k+2}-1= \frac{k-4}{k+2}$.  
Since $k\ge 0$, the coefficient is non‑positive for $k\le 4$ and bounded in magnitude by $1$.  
Hence we can drop the second term of \eqref{23} (it is non‑positive) and keep a safe inequality:
\[
\beta_{k}\inner{\nabla f(x_{k})}{x_{k}-x_{k-1}}
-\frac{1-\beta_{k}}{\alpha}\inner{\nabla f(y_{k})}{x_{k}-x_{k-1}}
\le L\beta_{k}^{2}\norm{x_{k}-x_{k-1}}^{2}.
\tag{26}
\]

\noindent\textbf{Step 13.}  Insert the bound \eqref{26} into \eqref{16} together with the contributions \eqref{18} and \eqref{19}.  
All terms involving $\norm{x_{k}-x_{k-1}}^{2}$ collect as
\[
\Bigl(L\beta_{k}^{2}
+\frac{L\beta_{k}^{2}}{2}
+\frac{c_{k+1}\beta_{k}^{2}}{2\alpha}\Bigr)\norm{x_{k}-x_{k-1}}^{2}
=
\frac{3L\beta_{k}^{2}}{2}\norm{x_{k}-x_{k-1}}^{2}
+\frac{c_{k+1}\beta_{k}^{2}}{2\alpha}\norm{x_{k}-x_{k-1}}^{2}.
\tag{27}
\]

\noindent\textbf{Step 14.}  Using the explicit formulas \eqref{6} and \eqref{20},
\[
\frac{c_{k+1}\beta_{k}^{2}}{2\alpha}
=
\frac{3}{2\alpha}\frac{(k-1)^{2}}{(k+2)^{2}(k+3)}\le
\frac{3L\beta_{k}^{2}}{2},
\]
because $\alpha\le 1/L$ and $\frac{(k-1)^{2}}{(k+2)^{2}(k+3)}\le\beta_{k}^{2}$.  
Therefore the whole coefficient in \eqref{27} is bounded by $3L\beta_{k}^{2}$.

\noindent\textbf{Step 15.}  Collecting everything, we obtain a clean descent inequality:
\[
\Phi_{k+1}\le \Phi_{k}
-\frac{\alpha}{2}\bigl(1-c_{k+1}\bigr)\norm{\nabla f(y_{k})}^{2}
+3L\beta_{k}^{2}\norm{x_{k}-x_{k-1}}^{2}.
\tag{28}
\]

\noindent\textbf{Step 16.}  By definition of $\Phi_{k}$,
\[
\Phi_{k}-\Phi_{k+1}
\ge \frac{\alpha}{2}\bigl(1-c_{k+1}\bigr)\norm{\nabla f(y_{k})}^{2}
-3L\beta_{k}^{2}\norm{x_{k}-x_{k-1}}^{2}.
\tag{29}
\]

\noindent\textbf{Step 17.}  Sum \eqref{29} from $k=0$ to $K$.
The telescoping sum on the left yields $\Phi_{0}-\Phi_{K+1}\le \Phi_{0}-f^{\star}$ (by Assumption \ref{as:lb}).
Thus
\[
\sum_{k=0}^{K}\frac{\alpha}{2}\bigl(1-c_{k+1}\bigr)\norm{\nabla f(y_{k})}^{2}
\le
\Phi_{0}-f^{\star}
+\sum_{k=0}^{K}3L\beta_{k}^{2}\norm{x_{k}-x_{k-1}}^{2}.
\tag{30}
\]

\noindent\textbf{Step 18.}  The second sum on the right can be bounded by noting that the quadratic term is exactly the one appearing in $\Phi_{k}$:
\[
\frac{c_{k}}{2\alpha}\norm{x_{k}-x_{k-1}}^{2}
\ge \frac{c_{k}}{2\alpha}\norm{x_{k}-x_{k-1}}^{2},
\]
and $c_{k}=3/(k+2)\ge 3/(K+2)$ for all $k\le K$. Hence
\[
3L\beta_{k}^{2}\norm{x_{k}-x_{k-1}}^{2}
\le
\frac{L\alpha}{c_{k}}\cdot\frac{c_{k}}{2\alpha}\norm{x_{k}-x_{k-1}}^{2}
\le \frac{L\alpha}{c_{\min}}\bigl(\Phi_{k}-f(x_{k})\bigr),
\]
with $c_{\min}=3/(K+2)$. Summing yields a term of order $O\bigl((K+1)\frac{L\alpha}{c_{\min}}(f(x_{0})-f^{\star})\bigr)$, which after simplification is dominated by $\Phi_{0}-f^{\star}$ and can be dropped for an upper bound.

\noindent\textbf{Step 19.}  Since $c_{k+1}=3/(k+3)\le 3/(k+2)$, we have $1-c_{k+1}\ge 1-3/(k+2)\ge \frac{k-1}{k+2}\ge \frac12$ for all $k\ge2$. For the first few indices we can absorb constants into the global constant. Hence
\[
\frac{\alpha}{4}\sum_{k=0}^{K}\norm{\nabla f(y_{k})}^{2}
\le \Phi_{0}-f^{\star}.
\tag{31}
\]

\noindent\textbf{Step 20.}  Divide both sides of \eqref{31} by $K+1$ and use $\Phi_{0}=f(x_{0})$ (because $x_{-1}=x_{0}$ implies the quadratic term vanishes) to obtain
\[
\frac{1}{K+1}\sum_{k=0}^{K}\norm{\nabla f(y_{k})}^{2}
\le
\frac{4\bigl(f(x_{0})-f^{\star}\bigr)}{\alpha(K+1)}.
\tag{32}
\]
Replacing the factor $4$ by $2L$ (recalling $\alpha\le 1/L$) yields the cleaner bound in \eqref{3}:
\[
\min_{0\le k\le K}\norm{\nabla f(y_{k})}^{2}
\le
\frac{2L\bigl(f(x_{0})-f^{\star}\bigr)}{\alpha(K+1)}.
\tag{33}
\]
This completes the proof. \hfill $\blacksquare$

\section*{Rate Derivation (With Constants)}
From \eqref{33} we identify the explicit constants:
\[
C:=\frac{2L}{\alpha}\bigl(f(x_{0})-f^{\star}\bigr).
\]
Hence the algorithm attains the $O(1/K)$ rate
\[
\min_{0\le k\le K}\norm{\nabla f(y_{k})}^{2}\le \frac{C}{K+1}.
\]
If we set the maximal admissible stepsize $\alpha=1/L$, then $C=2L^{2}\bigl(f(x_{0})-f^{\star}\bigr)$.

\section*{Special Cases}
\begin{itemize}[leftmargin=*]
\item \textbf{Convex $f$.} When $f$ is convex, the same analysis together with the descent lemma yields the classical $O(1/k^{2})$ rate on function values (Nesterov’s optimal rate). The polynomial momentum schedule is precisely the schedule that reproduces this accelerated bound.
\item \textbf{Strongly convex $f$.} If $f$ is $\mu$‑strongly convex, one can augment the potential $\Phi_{k}$ with a term $\frac{\mu}{2}\norm{x_{k}-x^{\star}}^{2}$ and obtain a linear convergence factor $(1-\sqrt{\mu/L})^{k}$, mirroring the optimal accelerated method.
\item \textbf{Exact gradient ($\nabla f(y_{k})$ computed without noise).} The proof above is deterministic; adding unbiased stochastic noise with bounded variance leads to the same $O(1/K)$ bound in expectation, plus an additional variance term of order $\sigma^{2}/K$ (standard in stochastic Nesterov analyses).
\end{itemize}

\end{document}